\section{Background}
\label{sec:background}

\subsection{Headless Browsers and Chrome DevTools Protocol (CDP)}

A headless browser executes JavaScript, renders the DOM, and responds to automation commands without a graphical display. Google Chrome introduced native headless mode in 2017 and quickly became a widely-used choice for browser automation due to its comprehensive web standards support and large tooling ecosystem. Its primary drawback in CI contexts is resource intensity: Chrome's multi-process architecture spawns dedicated renderer, GPU, and network service processes, resulting in hundreds of megabytes of resident memory per instance \cite{reis2009isolating}.

The CDP is a WebSocket-based remote control interface that exposes browser internals -- DOM manipulation, JavaScript execution, navigation, and network interception -- to external automation tools \cite{cdp2024}. Because Lightpanda implements a CDP-compatible server, Playwright can connect to it using the same endpoint configuration as for Chrome, making the two engines nominally interchangeable at the protocol level.

\subsection{Playwright and E2E Testing in CI}

Playwright is an open-source browser automation framework by Microsoft that provides high-level APIs for simulating user interactions: clicks, form input, navigation, and DOM assertions \cite{playwright2024}. It connects to browsers via CDP and supports \emph{browser contexts} -- isolated browsing sessions analogous to incognito windows -- for test isolation. Modern E2E test suites rely heavily on these contexts to prevent state leakage between tests and ensure reproducibility \cite{playwright2024}.

In CI pipelines such as GitHub Actions, E2E tests run on short-lived cloud runners with fixed resource quotas. Total pipeline duration is dominated not only by test execution but by setup overhead (browser download and initialisation), which Chrome's large binary makes expensive on cold runners \cite{hilton2016usage}. A lightweight alternative that reduces both setup and execution cost would be directly beneficial, provided it preserves compatibility.

\subsection{Lightpanda's Published Benchmark}

Lightpanda's homepage claims 9$\times$ faster execution and 16$\times$ lower memory consumption compared to Chrome, based on a ``crawler benchmark'' that navigates 933 ``real web pages.'' Inspection of the benchmark repository\footnote{\url{https://github.com/lightpanda-io/demo/blob/main/BENCHMARKS.md}} reveals that these pages are instances of a purpose-built Lightpanda demo application (\texttt{demo-browser.lightpanda.io/amiibo/}): a trivial static HTML page consisting of a single image and a small number of internal links. This is not a representative real-world web application. The validity of performance claims derived from such a benchmark is therefore limited, and independent empirical evaluation against real web applications is necessary.
